\section{Conclusion}
We investigated how acute stress affects interaction in immersive VR across components with distinct cognitive and perceptual demands. Drawing on Multiple Resource Theory as a task-selection framework, we examined three controlled probes within a unified VR environment, target acquisition, reading aloud, and ToL task, under a VR adaptation of the Trier Social Stress Test, combining behavioural performance, speech features, and physiological responses with self-reported anxiety.

Our results do not support a uniform-degradation account of stress in VR. Predefined target acquisition task outcomes remained broadly stable, while an exploratory measure of terminal pointing stability suggested subtler changes in fine-grained endpoint control. Reading aloud showed elevated mean F0 with comparatively stable timing, indicating that constrained verbal production preserves coarse temporal structure while still reflecting prosodic arousal. ToL task showed shorter total duration and execution time after stress, without corresponding changes in first-move latency or move count, suggesting altered execution dynamics rather than improved planning. Together, these patterns indicate that stress in VR manifests as selective, task-dependent behavioural modulation rather than as a single global decrement.

These findings extend situational-impairment research to immersive environments and offer an MRT-informed task-probe framing for reasoning about stress-related SIID in complex VR tasks. They also carry concrete design implications: stress-aware VR systems should adapt support to the current task rather than apply a generic response, monitor interaction dynamics beyond overt errors, and consider that post-stress effects may persist into recovery. More broadly, our work points toward stress-aware immersive systems that account jointly for the user's stress state and the resource demands of the interaction at hand. This could be enabled by real-time physiological sensing (e.g., wearable HR/HRV monitoring), allowing systems to dynamically infer stress states and adapt interaction support accordingly. Importantly, such adaptation should consider not only stress levels, but also the specific interaction demands of the task. Beyond immediate design implications, our findings also point toward future research that examines individual differences in stress response and the behavioural traces of stress induction itself within VR.